\documentclass[11pt]{article}

\usepackage[margin=1in]{geometry}

\usepackage{amsmath}

\usepackage[T1]{fontenc}

\title{Mathematical Modeling and Numerical Simulation Supplement}

\author{Anonymous Research Plan Author}

\date{}

\begin{document}

\maketitle

\section{Q1 (final: v0)}

\noindent\textbf{Abstract—} A parsimonious ordinary differential equation framework modeling the coupled evolution of neutron star spin period and magnetic inclination angle under vacuum dipole braking, wind enhancement, inclination decay, and accretion-driven spin-up. The model isolates the dynamical signatures of isolated magnetic-dipole evolution versus binary recycling by encoding phase-dependent torque switching directly into the state derivatives, enabling population-level statistical comparison against observed braking indices and characteristic ages.

\subsection{Question and Scope}

\noindent\textbf{Question.} Can the joint distribution of braking index n and characteristic age tau\_c statistically distinguish isolated magnetic-dipole evolution from binary recycling populations, given specific torque model constraints?



First executable model: parsimonious ODE integration of single pulsar tracks for isolated magnetic-dipole spin-down with optional inclination decay and wind braking, embedded in a Monte Carlo population synthesis framework. Binary recycling is represented as an accretion-phase switch with an effective spin-up law. The model excludes glitches, magnetospheric state changes, variable moment of inertia, detailed companion ablation, and full binary evolution.

\subsection{Model Assumptions and Applicability}

\begin{itemize}
\item ASM-001: Stellar geometry, electromagnetic constants, and moment of inertia are absorbed into effective braking constants K\_dipole and K\_acc. If violated: Derived constants lose direct physical interpretability and may require explicit microphysical coupling.
\item ASM-002: Magnetic inclination evolves smoothly via exponential decay without discontinuous glitches or magnetospheric state transitions. If violated: Smooth ODE integration fails to capture abrupt angular momentum transfers, invalidating continuous derivative assumptions.
\item ASM-003: Accretion torque is active exclusively during t < t\_acc with nonzero Mdot; isolated dipole torque applies otherwise. If violated: Simultaneous application of opposing torques produces unphysical period trajectories outside the validated regime.
\end{itemize}

\subsection{Symbols and Units}

\begin{center}
\small
\begin{tabular}{|p{0.12\linewidth}|p{0.15\linewidth}|p{0.46\linewidth}|p{0.17\linewidth}|}
\hline
ID & Symbol & Meaning & Unit \\ \hline
S-P & $P$ & Spin period & s \\ \hline
S-alpha & $\alpha$ & Magnetic inclination angle & rad \\ \hline
S-t & $t$ & Integration time & s \\ \hline
S-P0 & $P_0$ & Initial spin period & s \\ \hline
S-alpha0 & $\alpha_0$ & Initial inclination angle & rad \\ \hline
S-B0 & $B_0$ & Surface magnetic field strength & T \\ \hline
S-Kdip & $K_{\mathrm{dipole}}$ & Vacuum dipole braking constant & $s_T^-2$ \\ \hline
S-eta & $\eta$ & Wind braking efficiency & dimensionless \\ \hline
S-tau & $\tau_\alpha$ & Inclination decay timescale & s \\ \hline
S-Mdot & $\dot{M}$ & Mass accretion rate & $kg_s^-1$ \\ \hline
S-Kacc & $K_{\mathrm{acc}}$ & Accretion spin-up constant & $s^2_kg^-1$ \\ \hline
S-tacc & $t_{\mathrm{acc}}$ & Accretion phase duration & s \\ \hline
S-Tmax & $T_{\mathrm{max}}$ & Maximum simulated age & s \\ \hline
S-Pdotmin & $P_{\mathrm{dot,min}}$ & Minimum detectable period derivative & dimensionless \\ \hline
S-Nmc & $N_{\mathrm{mc}}$ & Monte Carlo sample count & dimensionless \\ \hline
\end{tabular}
\end{center}

\subsection{Mechanistic Equations}

\begin{equation}\label{eq:Q1-EQ-001}
\frac{dP}{dt} = \begin{cases} -\frac{K_{\mathrm{acc}} \dot{M}}{\max(P, 0.01)} & \mathrm{if} t < t_{\mathrm{acc}} \land \dot{M} > 0 \\ +\frac{K_{\mathrm{dipole}} B_0^2 \sin^2(\alpha)(1+\eta)}{\max(P, 0.01)} & \mathrm{otherwise} \end{cases}
\end{equation}
\noindent\textbf{Q1-EQ-001 where} S-P: Spin period; S-t: Integration time; S-Kacc: Accretion spin-up constant; S-Mdot: Mass accretion rate; S-tacc: Accretion phase duration; S-Kdip: Vacuum dipole braking constant; S-B0: Surface magnetic field strength; S-alpha: Magnetic inclination angle; S-eta: Wind braking efficiency.

\begin{equation}\label{eq:Q1-EQ-002}
\frac{d\alpha}{dt} = -\frac{\alpha}{\tau_\alpha}
\end{equation}
\noindent\textbf{Q1-EQ-002 where} S-alpha: Magnetic inclination angle; S-t: Integration time; S-tau: Inclination decay timescale.

\begin{equation}\label{eq:Q1-EQ-003}
n(t) = \frac{P(t) \cdot \frac{d^2P}{dt^2}}{\left(\frac{dP}{dt}\right)^2}
\end{equation}
\noindent\textbf{Q1-EQ-003 where} S-P: Spin period.

\begin{equation}\label{eq:Q1-EQ-004}
\tau_c(t) = \frac{P(t)}{2 \cdot \frac{dP}{dt}}
\end{equation}
\noindent\textbf{Q1-EQ-004 where} S-P: Spin period.

\subsection{Initial Conditions, Boundary Conditions, and Constraints}

\paragraph{Initial Conditions}
\begin{itemize}
\item \{'state': 'S-P', 'value': 0.5\}
\item \{'state': 'S-alpha', 'value': 0.7853981633974483\}
\end{itemize}

\paragraph{Boundary Conditions}
\begin{itemize}
\item \{'condition': 'Integration terminates when t reaches T\_max or |dP/dt| falls below Pdot\_min', 'type': 'TERMINATION'\}
\end{itemize}

\paragraph{Objective and Constraints}
\begin{itemize}
\item \{'objective': 'Generate distinct P-Pdot-n trajectories for isolated and recycled populations to enable statistical discrimination.'\}
\item \{'constraint': 'Sign convention: positive dP/dt denotes isolated spin-down; negative dP/dt denotes accretion spin-up. Torque branches must not activate simultaneously.'\}
\end{itemize}

\subsection{Algorithm}

\noindent\textbf{Algorithm}
\paragraph{Input}
\begin{itemize}
\item Initial conditions (P0, alpha0)
\item Approved parameter set
\item Scenario configuration
\end{itemize}
\paragraph{Output}
\begin{itemize}
\item Time series arrays for P(t), alpha(t)
\item Derived diagnostics n(t), tau\_c(t)
\item Population histograms
\end{itemize}
\paragraph{Steps}
\begin{itemize}
\item Initialize state vector with P0 and alpha0
\item Select scenario parameter overrides
\item Evaluate derivative AST at current time t
\item Advance state using fixed-step integrator
\item Compute n(t) and tau\_c(t) from derivatives
\item Check termination criteria against T\_max and Pdot\_min
\item Repeat until termination or N\_mc tracks completed
\end{itemize}

\subsection{Parameters, Scenarios, and Numerical Settings}

\paragraph{Parameters}
\begin{itemize}
\item \{'symbol\_id': 'S-P0', 'value': 0.5\}
\item \{'symbol\_id': 'S-alpha0', 'value': 0.7853981633974483\}
\item \{'symbol\_id': 'S-B0', 'value': 100000000.0\}
\item \{'symbol\_id': 'S-Kdip', 'value': 9.765625e-32\}
\item \{'symbol\_id': 'S-eta', 'value': 0.1\}
\item \{'symbol\_id': 'S-tau', 'value': 315576000000000.0\}
\item \{'symbol\_id': 'S-Mdot', 'value': 630000000000000.0\}
\item \{'symbol\_id': 'S-Kacc', 'value': 6.3e-31\}
\item \{'symbol\_id': 'S-tacc', 'value': 315576000000000.0\}
\item \{'symbol\_id': 'S-Tmax', 'value': 3155760000000000.0\}
\item \{'symbol\_id': 'S-Pdotmin', 'value': 1e-22\}
\item \{'symbol\_id': 'S-Nmc', 'value': 5000.0\}
\end{itemize}

\paragraph{Scenarios}
\begin{itemize}
\item \{'scenario\_id': 'isolated\_dipole', 'parameter\_overrides': \{'K\_acc': 0.0, 'Mdot': 0.0, 't\_acc': 0.0\}\}
\item \{'scenario\_id': 'binary\_recycling', 'parameter\_overrides': \{'K\_acc': 5e-31\}\}
\end{itemize}

\noindent\textbf{Solver family.} Explicit Runge-Kutta 4th Order.

\subsection{Parameter Provenance and Applicability}

\noindent\textbf{Approved parameter-set identity.} 0f0174315925c8e5b143507cd77c961e1becdad36afebaa31e7f6df12d27946c.
\begin{itemize}
\item initial\_period (P0) \ensuremath{=} 0.5 s; role BOUNDARY\_CONDITION; status APPROVED\_USER\_INPUT; source: Birth and Evolution of Isolated Radio Pulsars; DOI 10.1086/501516; document Q1-FGK2006; locator: preprint; applicability: required\_condition\_1\ensuremath{=}For isolated or recycled neutron stars at the beginning of the simulated spin-down phase, with period measured in seconds and no glitch-induced period jump at the initial epoch.. uncertainty: not reported.
\item initial\_inclination (alpha0) \ensuremath{=} 0.7853981633974483 rad; role BOUNDARY\_CONDITION; status APPROVED\_MODEL\_ASSUMPTION; source: MODEL ASSUMPTION (not a measured or literature value); applicability: required\_condition\_1\ensuremath{=}For dipole torque models with approximately aligned magnetic and rotational axes represented by an angle in radians; treated as the starting value before inclination decay.. uncertainty: not reported.
\item initial\_surface\_magnetic\_field (B0) \ensuremath{=} 100000000.0 T; role BOUNDARY\_CONDITION; status APPROVED\_USER\_INPUT; source: Binary and Millisecond Pulsars; DOI 10.12942/lrr-2008-8; document Q1-LORIMER2008; locator: review\_article; applicability: required\_condition\_1\ensuremath{=}For isolated neutron stars with approximately dipolar surface fields inferred from timing data, expressed in tesla; not intended for magnetar-like or strongly multipolar fields.. uncertainty: not reported.
\item vacuum\_dipole\_braking\_constant (K\_dipole) \ensuremath{=} 9.765625e-32 s\_T\textasciicircum{}-2; role MODEL\_ASSUMPTION; status APPROVED\_MODEL\_ASSUMPTION; source: MODEL ASSUMPTION (not a measured or literature value); applicability: required\_condition\_1\ensuremath{=}For the chosen effective spin-down law with magnetic field in tesla and period in seconds, assuming fixed stellar geometry and no unresolved internal braking physics.. uncertainty: not reported.
\item wind\_braking\_efficiency (eta) \ensuremath{=} 0.1 dimensionless; role MODEL\_ASSUMPTION; status APPROVED\_MODEL\_ASSUMPTION; source: MODEL ASSUMPTION (not a measured or literature value); applicability: required\_condition\_1\ensuremath{=}For pulsars with magnetospheric wind torque included as a multiplicative enhancement of the dipole term; scenario switches may disable the wind contribution.. uncertainty: not reported.
\item inclination\_decay\_timescale (tau\_alpha) \ensuremath{=} 315576000000000.0 s; role MODEL\_ASSUMPTION; status APPROVED\_USER\_INPUT; source: Binary and Millisecond Pulsars; DOI 10.12942/lrr-2008-8; document Q1-LORIMER2008; locator: review\_article; applicability: required\_condition\_1\ensuremath{=}For isolated pulsars modeled with smooth inclination decay, expressed in seconds, excluding glitches and abrupt magnetospheric state changes.. uncertainty: not reported.
\item accretion\_rate (Mdot) \ensuremath{=} 630000000000000.0 kg\_s\textasciicircum{}-1; role SCENARIO\_INPUT; status APPROVED\_MODEL\_ASSUMPTION; source: MODEL ASSUMPTION (not a measured or literature value); applicability: required\_condition\_1\ensuremath{=}For recycled pulsars in a steady disk-accretion phase, expressed in kilograms per second, ignoring transient outbursts and companion ablation.. uncertainty: not reported.
\item accretion\_braking\_constant (K\_acc) \ensuremath{=} 6.3e-31 s\textasciicircum{}2\_kg\textasciicircum{}-1; role SCENARIO\_INPUT; status APPROVED\_MODEL\_ASSUMPTION; source: MODEL ASSUMPTION (not a measured or literature value); applicability: required\_condition\_1\ensuremath{=}For the effective accretion law using accretion rate in kilograms per second and dimensionless period derivative; active only during the accretion phase.. uncertainty: not reported.
\item accretion\_duration (t\_acc) \ensuremath{=} 315576000000000.0 s; role SCENARIO\_INPUT; status APPROVED\_MODEL\_ASSUMPTION; source: MODEL ASSUMPTION (not a measured or literature value); applicability: required\_condition\_1\ensuremath{=}For binary recycling tracks represented by a finite accretion interval followed by isolated spin-down, expressed in seconds.. uncertainty: not reported.
\item maximum\_simulated\_age (T\_max) \ensuremath{=} 3155760000000000.0 s; role BOUNDARY\_CONDITION; status APPROVED\_USER\_INPUT; source: Birth and Evolution of Isolated Radio Pulsars; DOI 10.1086/501516; document Q1-FGK2006; locator: preprint; applicability: required\_condition\_1\ensuremath{=}Numerical endpoint for population synthesis tracks; not a physical observable unless tied to survey selection or evolutionary assumptions.. uncertainty: not reported.
\item minimum\_detectable\_period\_derivative (Pdot\_min) \ensuremath{=} 1e-22 dimensionless; role BOUNDARY\_CONDITION; status APPROVED\_MODEL\_ASSUMPTION; source: MODEL ASSUMPTION (not a measured or literature value); applicability: required\_condition\_1\ensuremath{=}Survey-dependent detection boundary for timing observations; used as a stopping condition rather than a physical spin-down parameter.. uncertainty: not reported.
\item monte\_carlo\_sample\_count (N\_mc) \ensuremath{=} 5000.0 dimensionless; role MODEL\_ASSUMPTION; status APPROVED\_MODEL\_ASSUMPTION; source: MODEL ASSUMPTION (not a measured or literature value); applicability: required\_condition\_1\ensuremath{=}Numerical control for population synthesis; chosen for computational budget and statistical resolution rather than physical measurement.. uncertainty: not reported.
\end{itemize}

\subsection{Numerical Validation and Model-Internal Results}

\paragraph{Validation Plan}
\begin{itemize}
\item \{'method': 'Analytical limit check: verify dalpha/dt converges to zero as t approaches infinity', 'metric': 'Relative error < 1e-6'\}
\item \{'method': 'Scenario isolation test: confirm isolated\_dipole yields strictly positive dP/dt throughout', 'metric': 'Sign consistency across all t'\}
\item \{'method': 'Recycling transition test: verify dP/dt switches sign at t\_acc and returns to positive post-accretion', 'metric': 'Branch activation count matches scenario definition'\}
\end{itemize}

\begin{itemize}
\item Execution sim-ef92f26db7c14196aaab75759811335e; NUMERICAL\_SIMULATION; SIMULATED; NOT\_EMPIRICAL; relation CONSTRAINED; summary: Within the audited reduced ODE, the isolated track increases from P\ensuremath{=}0.500 s to 0.662 s, whereas the recycling track reaches P\_min\ensuremath{=}0.226 s during accretion and ends at 0.281 s after resumed dipole spin-down. This establishes model-internal directional separation, but does not by itself validate population-level n-tau\_c discrimination because the executable MathIR contains two deterministic tracks rather than the documented Monte Carlo outer loop.
\end{itemize}

\subsection{Hypothesis Iteration Lineage}

\begin{itemize}
\item v0: CONSTRAINED (Within the audited reduced ODE, the isolated track increases from P\ensuremath{=}0.500 s to 0.662 s, whereas the recycling track reaches P\_min\ensuremath{=}0.226 s during accretion and ends at 0.281 s after resumed dipole spin-down. This establishes model-internal directional separation, but does not by itself validate population-level n-tau\_c discrimination because the executable MathIR contains two deterministic tracks rather than the documented Monte Carlo outer loop.)
\end{itemize}

\subsection{Limitations, Non-Empirical Disclosure, and References}

\noindent All values in this section are model-internal numerical simulations (NUMERICAL\_SIMULATION; SIMULATED; NOT\_EMPIRICAL), not empirical observations.

\paragraph{Limitations}
\begin{itemize}
\item Excludes glitch events, magnetospheric state transitions, and variable moment of inertia.
\item Assumes fixed stellar geometry and absorbs unresolved electromagnetic factors into effective constants.
\item Observational selection effects and survey detection thresholds are not dynamically coupled to the simulation loop.
\end{itemize}

\paragraph{References}
\begin{itemize}
\item \{'citation': 'Feng, G., et al. (2006). Birth and Evolution of Isolated Radio Pulsars. ApJ, 646, 381. DOI: 10.1086/501516', 'source': 'Q1-FGK2006'\}
\item \{'citation': 'Lorimer, D. R. (2008). Binary and Millisecond Pulsars. Living Rev. Relativity, 11, 8. DOI: 10.12942/lrr-2008-8', 'source': 'Q1-LORIMER2008'\}
\end{itemize}

\section{Q2 (final: v0)}

\noindent\textbf{Abstract—} A reduced-order ordinary differential equation initial-value model simulates the thermal evolution of isolated massive neutron stars under magnetic-dipole spin-down. The system tracks core temperature, surface temperature, central density, spin frequency, and a phase-fraction proxy. Scenario-specific switches activate latent-heat release and superfluid-modified cooling proxies when spin-down compression crosses a critical density threshold, enabling direct comparison of standard cooling, phase-transition cooling, and superfluid-enhanced phase-transition tracks.

\subsection{Question and Scope}

\noindent\textbf{Question.} How does the surface temperature evolution differ between massive neutron stars undergoing spin-down induced core compression and those without, and can this distinguish the phase-transition mechanism from standard cooling?



Isolated massive neutron star thermal evolution under magnetic-dipole spin-down, with a reduced initial-value model for core temperature, surface temperature, spin frequency, and central density. The model excludes binary recycling, companion ablation, detailed crustal heat blanketing, resolved neutrino microphysics, full equation-of-state integration, and non-spherical structure. Phase transition and superfluidity are represented through explicit scenario selection and effective coefficients.

\subsection{Model Assumptions and Applicability}

\begin{itemize}
\item ASM-001: Neutron stars are treated as spherically symmetric, isolated systems with no binary accretion torque or magnetic field evolution. If violated: Spin-down torque and rotational compression proxies become inaccurate; observed trajectories would diverge from modeled tracks.
\item ASM-002: Detailed neutrino emissivity processes and heat-capacity microphysics are absorbed into a single effective cooling timescale and scenario-dependent coefficients. If violated: Temperature relaxation rates would misrepresent actual stellar cooling physics, invalidating comparative scenario analysis.
\item ASM-003: Moment of inertia variations are fully parameterized into the rotational compression coefficient rather than solved dynamically. If violated: Central density evolution would decouple from realistic spin-down-induced structural changes.
\end{itemize}

\subsection{Symbols and Units}

\begin{center}
\small
\begin{tabular}{|p{0.12\linewidth}|p{0.15\linewidth}|p{0.46\linewidth}|p{0.17\linewidth}|}
\hline
ID & Symbol & Meaning & Unit \\ \hline
SYM-001 & $\Omega(t)$ & Angular spin frequency & $rad_s^-1$ \\ \hline
SYM-002 & $\rho_c(t)$ & Central density & $kg_m^-3$ \\ \hline
SYM-003 & $T_{\mathrm{core}}(t)$ & Core temperature & K \\ \hline
SYM-004 & $x(t)$ & Phase fraction proxy & dimensionless \\ \hline
SYM-005 & $T_{\mathrm{surface}}(t)$ & Surface temperature & K \\ \hline
SYM-006 & $M_{\mathrm{NS}}$ & Gravitational mass of the neutron star & $solar_mass$ \\ \hline
SYM-007 & $T_{c0}$ & Initial core temperature & K \\ \hline
SYM-008 & $\Omega_0$ & Initial angular spin frequency & $rad_s^-1$ \\ \hline
SYM-009 & $\rho_{c0}$ & Initial central density & $kg_m^-3$ \\ \hline
SYM-010 & $\dot{P}_0$ & Initial spin period derivative & $s_s^-1$ \\ \hline
SYM-011 & $\rho_{\mathrm{crit}}$ & Critical density for phase transition & $kg_m^-3$ \\ \hline
SYM-012 & $L_{\mathrm{latent}}$ & Effective latent-heat temperature proxy & K \\ \hline
SYM-013 & $S$ & Scenario selector switch & dimensionless \\ \hline
SYM-014 & $n$ & Braking index & dimensionless \\ \hline
SYM-015 & $\tau_{\mathrm{cool}}$ & Effective cooling timescale & s \\ \hline
SYM-016 & $\alpha_{cs}$ & Core-to-surface temperature conversion coefficient & dimensionless \\ \hline
SYM-017 & $\kappa_{\mathrm{rot}}$ & Rotational compression coefficient & $kg_s^2_m^-3$ \\ \hline
\end{tabular}
\end{center}

\subsection{Mechanistic Equations}

\begin{equation}\label{eq:Q2-EQ-001}
d\Omega/dt = -(\dot{P}_0 / 2\pi) \cdot \Omega_0^{2-n} \cdot \Omega^n
\end{equation}
\noindent\textbf{Q2-EQ-001 where} SYM-001: Angular spin frequency; SYM-008: Initial angular spin frequency; SYM-010: Initial spin period derivative; SYM-014: Braking index.

\begin{equation}\label{eq:Q2-EQ-002}
d\rho_c/dt = -\kappa_{\mathrm{rot}} \cdot \Omega \cdot (d\Omega/dt) \cdot (M_{\mathrm{NS}}/1.42)
\end{equation}
\noindent\textbf{Q2-EQ-002 where} SYM-002: Central density; SYM-001: Angular spin frequency; SYM-017: Rotational compression coefficient; SYM-006: Gravitational mass of the neutron star.

\begin{equation}\label{eq:Q2-EQ-003}
dx/dt = \begin{cases} (1-x)/(0.01\tau_{\mathrm{cool}}) & S \geq 1 \land \rho_c \geq \rho_{\mathrm{crit}} \\ 0 & \mathrm{otherwise} \end{cases}
\end{equation}
\noindent\textbf{Q2-EQ-003 where} SYM-004: Phase fraction proxy; SYM-013: Scenario selector switch; SYM-002: Central density; SYM-011: Critical density for phase transition; SYM-015: Effective cooling timescale.

\begin{equation}\label{eq:Q2-EQ-004}
dT_{\mathrm{core}}/dt = -\gamma \cdot T_{\mathrm{core}}/\tau_{\mathrm{cool}} + L_{\mathrm{latent}} \cdot dx/dt,   \gamma = \begin{cases} 1 & S < 2 \\ 0.5 & S \geq 2 \end{cases}
\end{equation}
\noindent\textbf{Q2-EQ-004 where} SYM-003: Core temperature; SYM-013: Scenario selector switch; SYM-015: Effective cooling timescale; SYM-012: Effective latent-heat temperature proxy; SYM-004: Phase fraction proxy.

\begin{equation}\label{eq:Q2-EQ-005}
dT_{\mathrm{surface}}/dt = \alpha_{cs} \cdot dT_{\mathrm{core}}/dt
\end{equation}
\noindent\textbf{Q2-EQ-005 where} SYM-005: Surface temperature; SYM-016: Core-to-surface temperature conversion coefficient; SYM-003: Core temperature.

\subsection{Initial Conditions, Boundary Conditions, and Constraints}

\paragraph{Initial Conditions}
\begin{itemize}
\item \{'condition\_id': 'IC-001', 'state': 'Omega', 'value': 30.44178928\}
\item \{'condition\_id': 'IC-002', 'state': 'rho\_c', 'value': 4.5e+17\}
\item \{'condition\_id': 'IC-003', 'state': 'T\_core', 'value': 1000000000.0\}
\item \{'condition\_id': 'IC-004', 'state': 'phase\_fraction', 'value': 0.0\}
\item \{'condition\_id': 'IC-005', 'state': 'T\_surface', 'value': 1000000.0\}
\end{itemize}

\paragraph{Boundary Conditions}
\begin{itemize}
\item \{'condition\_id': 'BC-001', 'description': 'Surface temperature initialized as linear projection of core temperature via alpha\_cs.', 'symbol\_refs': ['SYM-005', 'SYM-003', 'SYM-016']\}
\item \{'condition\_id': 'BC-002', 'description': 'Central density seeded below transition threshold to permit spin-down crossing.', 'symbol\_refs': ['SYM-002', 'SYM-011']\}
\end{itemize}

\paragraph{Objective and Constraints}
\begin{itemize}
\item \{'constraint\_id': 'OBJ-001', 'description': 'Maximize physical fidelity within reduced-order framework while maintaining computational tractability for scenario comparison.'\}
\item \{'constraint\_id': 'OBJ-002', 'description': 'Ensure all temperatures, densities, and phase fractions remain non-negative throughout integration.'\}
\item \{'constraint\_id': 'OBJ-003', 'description': 'Limit output trajectory points to 2000 or fewer by enforcing appropriate maximum step size relative to total simulation horizon.'\}
\end{itemize}

\subsection{Algorithm}

\noindent\textbf{Algorithm}
\paragraph{Input}
\begin{itemize}
\item Parameter set values
\item Scenario selector override
\item Time span bounds
\end{itemize}
\paragraph{Output}
\begin{itemize}
\item State trajectories over time
\item Surface temperature vs age tracks
\end{itemize}
\paragraph{Steps}
\begin{itemize}
\item Initialize state vector from approved initial conditions and parameter values.
\item Evaluate scenario selector to determine active cooling proxy and latent-heat activation logic.
\item Integrate coupled ODE system using stiff adaptive solver over defined time span.
\item Apply conditional branching for phase fraction evolution based on central density threshold.
\item Compute surface temperature trajectory from core temperature derivative and conversion coefficient.
\item Export synchronized time-series arrays for post-processing and observational comparison.
\end{itemize}

\subsection{Parameters, Scenarios, and Numerical Settings}

\paragraph{Parameters}
\begin{itemize}
\item \{'parameter\_id': 'neutron\_star\_mass', 'mathir\_symbol': 'M\_NS', 'value': 1.42, 'unit': 'solar\_mass', 'role': 'SCENARIO\_INPUT'\}
\item \{'parameter\_id': 'initial\_core\_temperature', 'mathir\_symbol': 'T\_c0', 'value': 1000000000.0, 'unit': 'K', 'role': 'BOUNDARY\_CONDITION'\}
\item \{'parameter\_id': 'initial\_spin\_frequency', 'mathir\_symbol': 'Omega0', 'value': 30.44178928, 'unit': 'rad\_s\textasciicircum{}-1', 'role': 'SCENARIO\_INPUT'\}
\item \{'parameter\_id': 'initial\_central\_density', 'mathir\_symbol': 'rho\_c0', 'value': 4.5e+17, 'unit': 'kg\_m\textasciicircum{}-3', 'role': 'BOUNDARY\_CONDITION'\}
\item \{'parameter\_id': 'initial\_spin\_period\_derivative', 'mathir\_symbol': 'Pdot0', 'value': 9.7228e-13, 'unit': 's\_s\textasciicircum{}-1', 'role': 'SCENARIO\_INPUT'\}
\item \{'parameter\_id': 'phase\_transition\_critical\_density', 'mathir\_symbol': 'rho\_crit', 'value': 5.12e+17, 'unit': 'kg\_m\textasciicircum{}-3', 'role': 'MATERIAL\_PROPERTY'\}
\item \{'parameter\_id': 'effective\_latent\_heat\_temperature\_release', 'mathir\_symbol': 'L\_latent', 'value': 50000000.0, 'unit': 'K', 'role': 'MATERIAL\_PROPERTY'\}
\item \{'parameter\_id': 'scenario\_selector', 'mathir\_symbol': 'scenario', 'value': 1.0, 'unit': 'dimensionless', 'role': 'SCENARIO\_INPUT'\}
\item \{'parameter\_id': 'braking\_index', 'mathir\_symbol': 'n', 'value': 3.15, 'unit': 'dimensionless', 'role': 'MODEL\_ASSUMPTION'\}
\item \{'parameter\_id': 'effective\_cooling\_timescale', 'mathir\_symbol': 'tau\_cool', 'value': 3155760000000.0, 'unit': 's', 'role': 'MODEL\_ASSUMPTION'\}
\item \{'parameter\_id': 'core\_surface\_conversion\_coefficient', 'mathir\_symbol': 'alpha\_cs', 'value': 0.001, 'unit': 'dimensionless', 'role': 'BOUNDARY\_CONDITION'\}
\item \{'parameter\_id': 'rotational\_compression\_coefficient', 'mathir\_symbol': 'kappa\_rot', 'value': 150000000000000.0, 'unit': 'kg\_s\textasciicircum{}2\_m\textasciicircum{}-3', 'role': 'MODEL\_ASSUMPTION'\}
\end{itemize}

\paragraph{Scenarios}
\begin{itemize}
\item \{'scenario\_id': 'standard\_cooling', 'parameter\_overrides': \{'scenario': 0.0\}, 'description': 'Standard cooling without phase transition or superfluid modification.'\}
\item \{'scenario\_id': 'phase\_transition', 'parameter\_overrides': \{'scenario': 1.0\}, 'description': 'Phase transition cooling with latent heat release upon density threshold crossing.'\}
\item \{'scenario\_id': 'phase\_transition\_superfluid', 'parameter\_overrides': \{'scenario': 2.0\}, 'description': 'Phase transition cooling with modified neutrino emissivity proxy (reduced cooling factor).'\}
\end{itemize}

\noindent\textbf{Solver family.} Stiff ODE IVP integrator (e.g., LSODA or Radau).

\subsection{Parameter Provenance and Applicability}

\noindent\textbf{Approved parameter-set identity.} 8960cd825beacfd9059c043b2efe1da4859a4586f8ae0739d68d043e5e953cc1.
\begin{itemize}
\item neutron\_star\_mass (M\_NS) \ensuremath{=} 1.42 solar\_mass; role SCENARIO\_INPUT; status APPROVED\_USER\_INPUT; source: Signature of Deconfinement with Spin-down Compression in Cooling Hybrid Stars; DOI 10.1088/0004-637X/694/2/1019; document Q2-STEJNER2009; locator: local\_document; applicability: required\_condition\_1\ensuremath{=}Isolated massive neutron star; spherical symmetry; non-recycled spin-down history; no binary accretion torque.. uncertainty: not reported.
\item initial\_core\_temperature (T\_c0) \ensuremath{=} 1000000000.0 K; role BOUNDARY\_CONDITION; status APPROVED\_USER\_INPUT; source: Signature of Deconfinement with Spin-down Compression in Cooling Hybrid Stars; DOI 10.1088/0004-637X/694/2/1019; document Q2-STEJNER2009; locator: local\_document; applicability: required\_condition\_1\ensuremath{=}Young isolated neutron star at the beginning of the cooling track; model assumes initial thermal state before long-term cooling; no crustal decoupling resolved.. uncertainty: not reported.
\item initial\_spin\_frequency (Omega0) \ensuremath{=} 30.44178928 rad\_s\textasciicircum{}-1; role SCENARIO\_INPUT; status APPROVED\_USER\_INPUT; source: Controlled excerpt from On the magnetic field evolution of PSR J1640-4631; DOI 10.3847/1538-4357/aa68e6; document Q2-GAO2017-PDOT-EXCERPT; locator: user\_provided\_local\_document; applicability: required\_condition\_1\ensuremath{=}Birth spin state of an isolated neutron star; magnetic dipole spin-down; no binary recycling or accretion-driven spin-up.. uncertainty: not reported.
\item initial\_central\_density (rho\_c0) \ensuremath{=} 4.5e+17 kg\_m\textasciicircum{}-3; role BOUNDARY\_CONDITION; status APPROVED\_MODEL\_ASSUMPTION; source: MODEL ASSUMPTION (not a measured or literature value); applicability: required\_condition\_1\ensuremath{=}Adopted for the first executable parsimonious model; must be internally consistent with the selected mass and EOS proxy; no full EOS integration is resolved.. uncertainty: not reported.
\item initial\_spin\_period\_derivative (Pdot0) \ensuremath{=} 9.7228e-13 s\_s\textasciicircum{}-1; role SCENARIO\_INPUT; status APPROVED\_USER\_INPUT; source: Controlled excerpt from On the magnetic field evolution of PSR J1640-4631; DOI 10.3847/1538-4357/aa68e6; document Q2-GAO2017-PDOT-EXCERPT; locator: user\_provided\_local\_document; applicability: required\_condition\_1\ensuremath{=}Isolated pulsar timing measurement; magnetic dipole braking; no accretion torque; no significant magnetic field evolution resolved.. uncertainty: not reported.
\item phase\_transition\_critical\_density (rho\_crit) \ensuremath{=} 5.12e+17 kg\_m\textasciicircum{}-3; role MATERIAL\_PROPERTY; status APPROVED\_USER\_INPUT; source: Signature of Deconfinement with Spin-down Compression in Cooling Hybrid Stars; DOI 10.1088/0004-637X/694/2/1019; document Q2-STEJNER2009; locator: local\_document; applicability: required\_condition\_1\ensuremath{=}Beta-equilibrated dense matter composition; EOS-dependent transition; relevant temperature and pressure regime for neutron star interiors.. uncertainty: not reported.
\item effective\_latent\_heat\_temperature\_release (L\_latent) \ensuremath{=} 50000000.0 K; role MATERIAL\_PROPERTY; status APPROVED\_MODEL\_ASSUMPTION; source: MODEL ASSUMPTION (not a measured or literature value); applicability: required\_condition\_1\ensuremath{=}EOS-dependent phase transition; beta-equilibrated dense matter; conversion from microscopic latent heat to a temperature proxy is a model assumption; no resolved heat-capacity microphysics.. uncertainty: not reported.
\item scenario\_selector (scenario) \ensuremath{=} 1.0 dimensionless; role SCENARIO\_INPUT; status APPROVED\_MODEL\_ASSUMPTION; source: MODEL ASSUMPTION (not a measured or literature value); applicability: required\_condition\_1\ensuremath{=}First executable parsimonious model; scenario choice is explicit and not represented as a literature value.. uncertainty: not reported.
\item braking\_index (n) \ensuremath{=} 3.15 dimensionless; role MODEL\_ASSUMPTION; status APPROVED\_USER\_INPUT; source: Controlled excerpt from On the magnetic field evolution of PSR J1640-4631; DOI 10.3847/1538-4357/aa68e6; document Q2-GAO2017-PDOT-EXCERPT; locator: user\_provided\_local\_document; applicability: required\_condition\_1\ensuremath{=}Adopted spin-down law assumption; no resolved magnetic field evolution; no accretion torque.. uncertainty: not reported.
\item effective\_cooling\_timescale (tau\_cool) \ensuremath{=} 3155760000000.0 s; role MODEL\_ASSUMPTION; status APPROVED\_MODEL\_ASSUMPTION; source: MODEL ASSUMPTION (not a measured or literature value); applicability: required\_condition\_1\ensuremath{=}Parsimonious first executable model; detailed Urca processes, modified Urca processes, direct Urca thresholds, and superfluid pairing are not resolved separately.. uncertainty: not reported.
\item core\_surface\_conversion\_coefficient (alpha\_cs) \ensuremath{=} 0.001 dimensionless; role BOUNDARY\_CONDITION; status APPROVED\_MODEL\_ASSUMPTION; source: MODEL ASSUMPTION (not a measured or literature value); applicability: required\_condition\_1\ensuremath{=}Spherical symmetry; simplified atmosphere boundary; no resolved radiative transfer or crustal thermal blanketing.. uncertainty: not reported.
\item rotational\_compression\_coefficient (kappa\_rot) \ensuremath{=} 150000000000000.0 kg\_s\textasciicircum{}2\_m\textasciicircum{}-3; role MODEL\_ASSUMPTION; status APPROVED\_MODEL\_ASSUMPTION; source: MODEL ASSUMPTION (not a measured or literature value); applicability: required\_condition\_1\ensuremath{=}Small spin-down induced density changes; no full rotating EOS integration; moment of inertia response is absorbed into the proxy.. uncertainty: not reported.
\end{itemize}

\subsection{Numerical Validation and Model-Internal Results}

\paragraph{Validation Plan}
\begin{itemize}
\item \{'plan\_id': 'VAL-001', 'method': 'Cross-scenario consistency check: verify that changing scenario selector alters at least one derivative term or state initialization.'\}
\item \{'plan\_id': 'VAL-002', 'method': 'Boundary limit test: confirm surface temperature remains positive and bounded for all three approved scenarios.'\}
\item \{'plan\_id': 'VAL-003', 'method': 'Observational bracketing: compare simulated T\_s(t) tracks against published X-ray cooling curves for isolated massive pulsars.'\}
\end{itemize}

\begin{itemize}
\item Execution sim-05200df7f99b4a56b313fb266667735f; NUMERICAL\_SIMULATION; SIMULATED; NOT\_EMPIRICAL; relation SUPPORTED\_WITHIN\_MODEL; summary: The audited reduced ODE crosses rho\_crit\ensuremath{=}5.12e17 kg/m\textasciicircum{}3 at approximately 9.85e11 s in all spin-down cases. At 3*tau\_cool, the standard, phase-transition, and phase-transition-plus-superfluid scenarios yield surface temperatures of about 4.98e4 K, 5.32e4 K, and 2.36e5 K, respectively; phase\_fraction reaches 1 only in the two transition scenarios. The mechanism therefore produces distinguishable model-internal thermal tracks, conditional on the effective cooling, latent-heat, and compression proxies.
\end{itemize}

\subsection{Hypothesis Iteration Lineage}

\begin{itemize}
\item v0: SUPPORTED\_WITHIN\_MODEL (The audited reduced ODE crosses rho\_crit\ensuremath{=}5.12e17 kg/m\textasciicircum{}3 at approximately 9.85e11 s in all spin-down cases. At 3*tau\_cool, the standard, phase-transition, and phase-transition-plus-superfluid scenarios yield surface temperatures of about 4.98e4 K, 5.32e4 K, and 2.36e5 K, respectively; phase\_fraction reaches 1 only in the two transition scenarios. The mechanism therefore produces distinguishable model-internal thermal tracks, conditional on the effective cooling, latent-heat, and compression proxies.)
\end{itemize}

\subsection{Limitations, Non-Empirical Disclosure, and References}

\noindent All values in this section are model-internal numerical simulations (NUMERICAL\_SIMULATION; SIMULATED; NOT\_EMPIRICAL), not empirical observations.

\paragraph{Limitations}
\begin{itemize}
\item Highly sensitive to the choice of EOS and phase transition parameters; observational uncertainties in distance and radius affect temperature inference.
\item Assumes spherical symmetry and ignores binary recycling, companion-ablation evolution, and detailed crustal heat blanketing.
\item No resolved neutrino microphysics or full equation-of-state integration; relies on effective proxy coefficients for first executable run.
\end{itemize}

\paragraph{References}
\begin{itemize}
\item \{'ref\_id': 'REF-001', 'citation': 'Stejner et al. (2009), Signature of Deconfinement with Spin-down Compression in Cooling Hybrid Stars, ApJ, 694, 1019.'\}
\item \{'ref\_id': 'REF-002', 'citation': 'Gao et al. (2017), On the magnetic field evolution of PSR J1640-4631, ApJ, 840, 12.'\}
\end{itemize}

\end{document}
